\ncuAbstractPage{\ncuCNtitle}{中文摘要}{摘\quad 要}
時間變化與類別不平衡常同時出現在分類問題中：歷史資料樣本較多，卻未必符合近期情況；近期資料較接近預測時點，少數類樣本卻較少。本研究以具有時間資訊的不平衡分類為對象，比較以歷史（Old）與近期（New）資料建立之模型的運用方式，並評估模型集成與特徵選取的影響。

實證使用 1999--2018 年美國企業的公開年度財務資料，以破產公司的最後觀測年度重建下一年度破產事件標籤，正類僅占全部紀錄的 0.77\%。實驗一在相同年份配置下比較四種分類器的 Old、New 與合併重訓（Retrain），再以 XGBoost 建立 Old／New 六模型池，比較靜態集成、動態分類器選擇與動態集成選擇；實驗二在採樣前加入四種特徵選取方法與兩種保留數。所有成員選擇與分類門檻均由驗證資料決定。

結果顯示，近期資料的優勢依分類器而異。在固定模型池中，只整合近期模型的 New3 平均 AUC 最高（0.8601），依驗證資料選出兩個成員的靜態集成 F1-score 最高（0.1795）；等權整合全部六模型與動態選擇則未帶來改善。特徵選取沒有在多數方法與指標上一致改善。調參後的單一分類器與集成表現相當，因此本研究的證據支持審慎選擇集成成員，而非主張集成必然優於單一模型。改用公司最終狀態標籤時，方法的排序隨之改變，顯示標籤定義會影響結論。

\noindent\textbf{關鍵字：}類別不平衡、時間資料、模型集成、動態選擇、特徵選取、破產預測
